%% CPP 2027 Submission — 0pat Exercise XXVI: Zero Reflection Symmetry (First Analytic Material of the Bridge)
\documentclass[sigplan,10pt,anonymous,review]{acmart}
\settopmatter{printfolios=true,printacmref=false}

\AtBeginDocument{%
  \providecommand\BibTeX{{\normalfont B\scshape ib\TeX}}}

\settopmatter{printacmref=false}
\renewcommand\footnotetextcopyrightpermission[1]{}
\pagestyle{plain}

\usepackage{microtype}
\usepackage{booktabs}
\usepackage{amsmath}
% XeTeX (tectonic): acmart loads newtxmath/libertine which already defines \Bbbk;
% reset it so amssymb's definition is accepted.
\let\Bbbk\relax
\usepackage{amssymb}

\begin{document}

\title{Zero Reflection Symmetry: The First Analytic Material of the
Riemann Bridge (Exercise XXVI)}

\author{Anonymous}
\affiliation{%
  \institution{Anonymous}
}
\email{anonymous@example.org}

\begin{abstract}
We report the first analytic theorem of the Riemann-direction chain: the
zero reflection symmetry of the Riemann zeta function, proved from the
classical functional equation in Lean 4 / mathlib with zero
\texttt{sorry}. The verification map of the Riemann hypothesis has two
completed banks --- the Euler-product bank on $\mathrm{Re}(s) \ge 1$
(zero-free region) and the geometric bank (critical line $\leftrightarrow$
critical circle) --- and one open span: the analytic control of zeros in
the critical strip $0 < \mathrm{Re}(s) < 1$. The functional equation
$\zeta(1-s) = 2(2\pi)^{-s}\Gamma(s)\cos(\pi s/2)\zeta(s)$, present in
mathlib as \texttt{riemannZeta\_one\_sub}, is the only known analytic
material that connects the two banks. This exercise formalizes its first
corollary: every nontrivial zero reflects to a zero
($\zeta(s) = 0 \Rightarrow \zeta(1-s) = 0$), the reflection stays in the
strip, and an off-line zero ($\mathrm{Re}(s) \ne 1/2$) therefore occurs
in pairs. Three theorems, zero errors, mathlib v4.32.2.
\end{abstract}

\keywords{Lean, formalization, Riemann zeta function, functional equation, zero symmetry, 0pat}

\maketitle

\section{The position of this exercise}

The Riemann-direction chain is a verification map with two completed
banks and one open span. The left bank is analytic: the Euler product
converges on $\mathrm{Re}(s) > 1$ and $\zeta$ has no zeros on
$\mathrm{Re}(s) \ge 1$ (Exercise C025). The right bank is geometric:
the critical line is the critical circle under inversion, and the
hypothetical zero positions and the critical circle are the same object
(Exercises C019--C022). The open span is the critical strip
$0 < \mathrm{Re}(s) < 1$: no theorem of the chain controls the positions
of the zeros of $\zeta$ there.

The only known analytic material that spans the two banks is the
functional equation. Mathlib states it in classical form
(\texttt{riemannZeta\_one\_sub}): for $s \ne -n$ ($n \in \mathbb{N}$)
and $s \ne 1$,

\[
\zeta(1-s) = 2(2\pi)^{-s}\,\Gamma(s)\,\cos(\pi s / 2)\,\zeta(s).
\]

This exercise takes the first corollary of this equation: zero
reflection symmetry. It is the first theorem of the chain that genuinely
uses analysis --- the functional equation --- rather than algebra or
geometry.

\section{The theorems}

Three theorems, all machine-checked with zero \texttt{sorry}.

\paragraph{Reflection (zeta\_zero\_reflection).} If $\zeta(s) = 0$ and
$s$ is not a trivial zero position ($s \ne -n$ for all $n$) and not the
pole ($s \ne 1$), then $\zeta(1-s) = 0$. The proof is one rewriting
step: substitute the functional equation and multiply the zero by the
coefficient. The coefficient need not be nonzero --- zero times any
coefficient is zero --- so no analytic nonvanishing is required.

\paragraph{The reflection stays in the strip (zeta\_zero\_strip\_reflection).}
If $0 < \mathrm{Re}(s) < 1$, then $0 < \mathrm{Re}(1-s) < 1$ as well.
Together with reflection: a zero in the strip reflects to a zero in the
strip. Zeros of $\zeta$ are therefore symmetric about the critical line
$\mathrm{Re}(s) = 1/2$ within the strip.

\paragraph{Off-line zeros come in pairs (zeta\_zero\_offline\_pair).}
If $\mathrm{Re}(s) \ne 1/2$ (the zero is off the critical line), then
$1-s \ne s$, and the reflection is a \emph{different} zero of the strip.
The proof uses the real-part computation: $1-s = s$ would force
$\mathrm{Re}(s) = 1/2$. This is the first structural constraint on
off-line zeros that the chain obtains: if an off-line zero exists, it is
accompanied by a second, distinct off-line zero. The Riemann hypothesis
asserts that no such pair exists; this theorem asserts only that pairs
are the only possible form of violation.

\section{Honest boundary and position in the map}

As in every exercise of the series: the content is classical mathematics
(``novelty\_status: KNOWN'' --- the reflection symmetry is a standard
corollary of the functional equation, known since Riemann), the proof is
machine-checked, and no claim is made of proving the Riemann hypothesis.
The value of the exercise is positional: it is the first analytic node of
the map's open span, it connects the two completed banks through the
functional equation, and it converts the question ``is there an off-line
zero?'' into the more structured question ``is there an off-line zero
pair?''. The bridge now has its first material; the span remains open.

\section{Formalization}

The proof is a single Lean file (\texttt{proof.lean}, 62 lines),
importing only \texttt{Mathlib.NumberTheory.LSeries.RiemannZeta}.
Compiled with Lean 4.32.2 / mathlib v4.32.2, zero errors, zero
\texttt{sorry}. The functional equation
\texttt{riemannZeta\_one\_sub} is mathlib's classical statement;
\texttt{Complex.sub\_re} and \texttt{linarith} handle the strip
inequalities.

\begin{center}
\begin{tabular}{lll}
\toprule
Theorem & Statement & Proof core \\
\midrule
\texttt{zeta\_zero\_reflection} & $\zeta(s)=0 \Rightarrow \zeta(1-s)=0$
  & \texttt{rw [riemannZeta\_one\_sub, hz]; ring} \\
\texttt{zeta\_zero\_strip\_reflection} & strip-preserving reflection
  & \texttt{Complex.sub\_re; linarith} \\
\texttt{zeta\_zero\_offline\_pair} & $\mathrm{Re}(s) \ne 1/2 \Rightarrow 1-s \ne s$ and a zero
  & real-part computation + reflection \\
\bottomrule
\end{tabular}
\end{center}

\section{Conclusion}

The first analytic material of the Riemann bridge is in place: zero
reflection symmetry, proved from the functional equation, machine-checked
in Lean. The map now reads: left bank (Euler product, zero-free region),
first span material (reflection symmetry, off-line pairs), right bank
(critical line $\leftrightarrow$ critical circle). The span remains open
--- the reflection symmetry permits off-line pairs and the Riemann
hypothesis denies them --- but the question it poses is now exactly
structured: can the analysis show that no off-line zero pair exists?

\bibliographystyle{ACM-Reference-Format}
\bibliography{refs}

\end{document}
